%% front-notation.tex
%% SAMBP thesis notation — loaded by \addNotationDefinitions / \printNotation
%% Format: \nomenclature{symbol}{description}

%% ── System quantities ────────────────────────────────────────
\nomenclature{$k_{\mathrm{ibr}}$}{IBR current ratio $= I_{\mathrm{IBR}}/I_{\mathrm{rated}}$, dimensionless}
\nomenclature{$I_{\mathrm{IBR}}$}{IBR terminal current magnitude (pu)}
\nomenclature{$I_{\mathrm{rated}}$}{IBR rated current (pu)}
\nomenclature{$I_{\mathrm{lim}}$}{IBR current limiter ceiling, 1.1--2.0\,pu}
\nomenclature{$\underline{I}_A,\,\underline{I}_B$}{Phasor currents at relay end A / end B}
\nomenclature{$V_1,\,V_2,\,V_0$}{Positive-, negative-, zero-sequence voltages}
\nomenclature{$I_1,\,I_2,\,I_0$}{Positive-, negative-, zero-sequence currents}
\nomenclature{$Z_{\mathrm{line}}$}{Positive-sequence line impedance ($\Omega$)}
\nomenclature{$Z_F$}{Fault impedance ($\Omega$)}
\nomenclature{$Z_{\mathrm{app}}$}{Apparent impedance seen by distance relay ($\Omega$)}
\nomenclature{$\alpha$}{Fault location parameter, $0\le\alpha\le1$}
\nomenclature{$\hat\alpha$}{EKF estimate of fault location parameter}
\nomenclature{$\hat k_{\mathrm{ibr}}$}{EKF estimate of IBR current ratio}
\nomenclature{$\hat Z_{\mathrm{line}}$}{EKF estimate of line impedance}

%% ── Differential protection ──────────────────────────────────
\nomenclature{$I_{\mathrm{diff}}$}{Differential operating current $=|\underline{I}_A+\underline{I}_B|$}
\nomenclature{$I_{\mathrm{bias}}$}{Bias (restraint) current $= (|\underline{I}_A|+|\underline{I}_B|)/2$}
\nomenclature{$k_m$}{Differential relay operating slope (fixed), 0.2--0.5}
\nomenclature{$k_m^*$}{Adaptive differential slope (function of $k_{\mathrm{ibr}}$)}
\nomenclature{$I_{\mathrm{min}}$}{Differential relay minimum-operating current (pu)}
\nomenclature{$k_{\mathrm{TDCS}}$}{Transient Differential Correlation Score (0--1)}

%% ── Distance protection ──────────────────────────────────────
\nomenclature{$Z_1$}{Zone-1 reach impedance ($\Omega$)}
\nomenclature{$k_r$}{IBR reach-correction factor}
\nomenclature{$\beta$}{IBR reactive current injection angle (rad)}

%% ── Machine learning ─────────────────────────────────────────
\nomenclature{$\mathcal{L}$}{Total PINN loss function}
\nomenclature{$\mathcal{L}_{\mathrm{data}}$}{Supervised cross-entropy loss}
\nomenclature{$\mathcal{L}_{C1},\,\mathcal{L}_{C2},\,\mathcal{L}_{C3}$}{Physics-constraint penalty terms}
\nomenclature{$\lambda_1,\,\lambda_2,\,\lambda_3$}{Penalty weights (10, 5, 2)}
\nomenclature{$S_k$}{CUSUM statistic at decision cycle $k$}
\nomenclature{$\theta$}{CUSUM alarm threshold (= 5)}
\nomenclature{$\delta$}{CUSUM slack (allowable drift per cycle)}

%% ── Kalman filter ────────────────────────────────────────────
\nomenclature{$\mathbf{x}_k$}{EKF state vector at time step $k$}
\nomenclature{$\mathbf{Q}$}{EKF process noise covariance matrix}
\nomenclature{$\mathbf{R}$}{EKF measurement noise covariance matrix}
\nomenclature{$\mathbf{H}_k$}{EKF measurement Jacobian}
\nomenclature{$\mathbf{P}_k$}{EKF error covariance matrix}
\nomenclature{$\mathbf{K}_k$}{Kalman gain matrix}

%% ── Islanding / ROCOF ────────────────────────────────────────
\nomenclature{$\mathrm{d}f/\mathrm{d}t$}{Rate of change of frequency (Hz/s)}
\nomenclature{$\gamma_1$}{ROCOF threshold (0.5\,Hz/s)}
\nomenclature{$\gamma_2$}{Phase deviation threshold ($2^\circ$)}
\nomenclature{$\gamma_3$}{Negative-sequence voltage ratio threshold (0.02)}

%% ── Optimisation ─────────────────────────────────────────────
\nomenclature{$t_{ij}$}{Operating time of relay $R_i$ for fault on feeder $j$}
\nomenclature{$M_{ij}$}{Plug-setting multiplier of relay $R_i$ for fault $j$}
\nomenclature{$\mathrm{TMS}_i$}{Time-multiplier setting of relay $R_i$}
\nomenclature{$\Delta_{\mathrm{CTI}}$}{Coordination time interval, $\ge 0.2$\,s}

%% ── Constants and operators ──────────────────────────────────
\nomenclature{$\jj$}{Imaginary unit $\sqrt{-1}$}
\nomenclature{$\mathrm{pu}$}{Per-unit system normalisation}
\nomenclature{$\underline{\,\cdot\,}$}{Complex phasor quantity}
\nomenclature{$|\,\cdot\,|$}{Magnitude of complex number}
\nomenclature{$\angle(\,\cdot\,)$}{Argument (phase angle) of complex number}
